\chapter{누가 더 빨리 계산할까}

6장에서 물 붓기를 손으로 해 봤지요.
친구가 네 명이고 화살표가 다섯 개였는데도 꽤 번거로웠어요.
진짜 지도에는 화살표가 수만, 수십만 개 있어요.
컴퓨터가 하더라도 화살표가 많을수록 시간이 더 걸려요.

\section{편지 나르기}

우체부 두 명이 있다고 해 봐요.
한 사람의 가방에는 편지가 1만 5천 통쯤, 다른 사람의 가방에는 13만 7천 통쯤 들어 있어요.
둘이 같은 속도로 일한다면 누가 먼저 끝날까요?
당연히 가방이 가벼운 쪽이에요.
가방의 무게가 아홉 배 차이 나면 걸리는 시간도 아홉 배쯤 차이 나겠지요.

엔도스랭크와 AWP가 바로 그래요.
8장에서 본 것처럼, 같은 지갑 5,521개로 지도를 그리면
엔도스랭크의 화살표는 14,727개, AWP의 화살표는 137,087개예요.

\section{실제로 재 보았어요}

같은 컴퓨터, 같은 프로그램, 같은 놀이 규칙으로 두 점수를 계산하고 시간을 쟀어요.
처음 한 번은 몸풀기로 재지 않고, 그다음 다섯 번을 재서 평균을 냈어요.

\begin{figure}[h]
\centering
\begin{tikzpicture}
\begin{axis}[
  xbar, bar width=11pt, width=0.92\linewidth, height=4.6cm,
  xmin=0, xmax=5.6, xlabel={계산에 걸린 시간 (초)},
  symbolic y coords={S, C, A, E}, ytick=data,
  yticklabels={S-PR, C-PR, AWP, 엔도스랭크},
  nodes near coords, nodes near coords style={font=\scriptsize},
  every node near coord/.append style={/pgf/number format/fixed, /pgf/number format/precision=3},
  tick label style={font=\small}, label style={font=\small}, axis y line*=left,
  enlarge y limits=0.2,
]
\addplot[fill=inkblue!70, draw=inkblue] coordinates {(4.156,S) (4.760,C) (4.105,A) (0.401,E)};
\end{axis}
\end{tikzpicture}
\caption{지갑 5,521개로 점수를 한 번 계산하는 데 걸린 평균 시간.}
\end{figure}

엔도스랭크는 0.401초, AWP는 4.105초가 걸렸어요.
엔도스랭크가 약 \textbf{10.2배} 빨랐어요.
화살표 수의 차이(약 9.3배)와 거의 같은 비율이에요.
기억하는 데 쓴 메모리도 4.6메가바이트와 38.2메가바이트로 비슷한 차이가 났어요.

시간의 대부분은 물을 붓는 데가 아니라, 계산을 시작하기 전에 화살표 목록을 표로 정리하는 데 들었어요.
제가 따로 재 보니 그 정리가 전체 시간의 94\%쯤이었어요.
그 정리에 걸리는 시간은 화살표 수에 비례하니까, 결국 화살표가 적은 쪽이 빠른 거예요.

\section{정직하게 덧붙일 것}

이 결과에서 말할 수 있는 것과 말할 수 없는 것을 나누어 둘게요.

\begin{itemize}
\item \textbf{말할 수 있는 것}: 같은 지갑들을 놓고 보면, 엔도스랭크는 화살표가 적어서 AWP보다 훨씬 빨리 계산돼요.
\item \textbf{말할 수 없는 것}: 지갑이 훨씬 많아져도 늘 이만큼 빠를지는 확인하지 못했어요.
      지갑 수를 10만 개 가까이 늘려서도 재 보았지만, 새로 더한 지갑들에는 허락이나 송금 화살표를 따로 가져오지 못했어요.
      그래서 지도 자체는 커지지 않았어요.
\item 두 지도를 섞는 C-PR은 화살표를 둘 다 쓰니 AWP보다 빠를 수 없어요. 실제로 4.76초가 걸렸어요.
      빠르다는 장점은 엔도스랭크에만 해당해요.
\end{itemize}

\begin{deeper}[진짜 숫자]
\begin{center}
\small
\begin{tabular}{@{}lrrrr@{}}
\toprule
& 시간(초) & 반복 수 & 점 & 화살표 \\
\midrule
엔도스랭크 & 0.401 & 37 & 6,442 & 14,727 \\
AWP & 4.105 & 100 & 30,791 & 137,087 \\
C-PR ($\lambda=0.5$) & 4.760 & 76 & 31,242 & 151,814 \\
S-PR & 4.156 & 145 & 30,791 & 137,087 \\
\bottomrule
\end{tabular}
\end{center}
시간은 몸풀기 1회 뒤 5회 잰 평균이에요. 지갑 수를 1,000개, 2,000개, 4,000개로 줄여 가며 잰 결과와,
확장 지갑 풀(10,000개, 50,000개, 99,080개)로 잰 결과에서도 엔도스랭크가 8.1배에서 12.2배 빨랐어요.
확장 풀의 새 지갑들은 화살표가 없어서 지도에 점을 더하지 못했어요. 그래서 99,080개 단계의 지도는 5,521개 단계와 같아요.
\end{deeper}

\begin{learned}
\begin{itemize}
\item 계산 시간은 화살표 수에 거의 비례해요.
\item 엔도스랭크는 화살표가 약 9배 적어서 AWP보다 약 10배 빨랐어요.
\item 지도가 훨씬 커져도 그런지는 아직 확인하지 못했어요.
\end{itemize}
\end{learned}
